\documentclass[11pt]{article}

\usepackage[utf8]{inputenc}
\usepackage{amsmath, amssymb, amsthm}
\usepackage{graphicx}
\usepackage{hyperref}
\usepackage{booktabs}
\usepackage{geometry}
\geometry{margin=1in}

\title{StochasTech: Differentiable Heston Calibration for Value-at-Risk Forecasting}
\author{Amine Cheikhrouhou}
\date{\today}

\begin{document}
\maketitle

\begin{abstract}
We present a simulation-based framework for financial risk analysis that
calibrates the Heston stochastic-volatility model directly to historical equity
returns by backpropagating through the SDE solver. The differentiable
full-truncation Euler scheme is wired to PyTorch autograd, allowing the five
Heston parameters $(\kappa, \theta, \xi, \rho, v_0)$ plus drift $\mu$ to be
fit by Adam against an energy-distance loss against observed log-returns.
Out-of-sample 1-day Value-at-Risk forecasts are evaluated on SPY, AAPL, and
MSFT against a closed-form GBM-MLE baseline and a historical-simulation
benchmark. Coverage is measured by the Kupiec proportion-of-failures and
Christoffersen independence/conditional-coverage tests. The contribution is
the combination of a simple, robust BPTT calibrator (no stochastic adjoint
required at this parameter scale) with an energy-distance loss that is
forgiving of model misspecification and produces stable parameter
trajectories under rolling-window calibration.
\end{abstract}

\section{Introduction}
\label{sec:intro}

Stochastic-volatility models are a workhorse of derivatives pricing and
quantitative risk management. The Heston model \cite{heston1993} introduced a
mean-reverting CIR process for the instantaneous variance, correlated with the
return-driving Brownian motion to capture the empirically observed leverage
effect. Its closed-form characteristic function makes it tractable for
European-option pricing, but the calibration problem against historical return
series remains hard: the variance is latent, and the simulator's likelihood is
not available in closed form.

Recent advances in differentiable programming \cite{li2020adjointsde,
chen2018neuralode} make it practical to differentiate through SDE solvers,
turning calibration into a gradient-based optimization problem. We exploit
this to fit Heston parameters by backpropagating through the full-truncation
Euler simulator \cite{lord2010biased} of the joint $(\log S_t, v_t)$ system,
using the energy distance \cite{szekely2013energy} between simulated and
observed return distributions as the loss.

We then evaluate the calibrated model in the standard VaR-backtesting
framework of Kupiec \cite{kupiec1995} and Christoffersen
\cite{christoffersen1998}, comparing against (i) closed-form GBM maximum
likelihood and (ii) historical simulation.

\textbf{Contributions.}
\begin{enumerate}
\item A self-contained, $\sim$100-line PyTorch implementation of a
  differentiable full-truncation Heston simulator that supports backprop with
  zero modification to the forward path beyond an $\varepsilon$-floor inside
  $\sqrt{v^+}$ for autograd safety.
\item A demonstration that for Heston's small parameter dimension, vanilla
  backprop-through-time is competitive with the stochastic adjoint
  \cite{li2020adjointsde} in wall-clock time and gives substantially simpler
  numerics.
\item Out-of-sample VaR coverage on three equity tickers, with rolling-window
  refit and standard backtesting tests.
\end{enumerate}

\section{Mathematical Model}
\label{sec:model}

Let $S_t$ denote the price of a single equity and $v_t$ its instantaneous
variance. The Heston system is
\begin{align}
  dS_t &= \mu S_t\, dt + \sqrt{v_t}\, S_t\, dW_t^S, \label{eq:heston-S}\\
  dv_t &= \kappa(\theta - v_t)\, dt + \xi \sqrt{v_t}\, dW_t^v, \label{eq:heston-v}\\
  d\langle W^S, W^v\rangle_t &= \rho\, dt. \label{eq:heston-rho}
\end{align}
Applying It\^o's lemma to $\log S_t$ gives the simulation-ready form
\begin{equation}
  d(\log S_t) = (\mu - \tfrac{1}{2} v_t)\, dt + \sqrt{v_t}\, dW_t^S.
  \label{eq:heston-logS}
\end{equation}
The variance process is a Cox--Ingersoll--Ross diffusion. The Feller
condition $2\kappa\theta \ge \xi^2$ ensures $v_t > 0$ a.s. in the continuous
SDE; discrete schemes can violate it, motivating the full-truncation
discretization below.

\section{Method: Differentiable Calibration}
\label{sec:method}

\subsection{Full-truncation Euler simulator}

For step size $\Delta t$ and grid $t_i = i \Delta t$, we discretize
\eqref{eq:heston-v}--\eqref{eq:heston-logS} as
\begin{align}
  v_{i+1} &= v_i + \kappa(\theta - v_i^+)\, \Delta t + \xi \sqrt{v_i^+}\, \sqrt{\Delta t}\, Z_i^v, \\
  \log S_{i+1} &= \log S_i + (\mu - \tfrac{1}{2} v_i^+)\, \Delta t
                + \sqrt{v_i^+}\, \sqrt{\Delta t}\, Z_i^S, \\
  Z_i^S &= \rho Z_i^v + \sqrt{1 - \rho^2}\, Z_i^\perp,
\end{align}
where $v_i^+ = \max(v_i, 0)$, and $Z_i^v, Z_i^\perp$ are independent standard
normals. The truncation only affects how $v_i$ is consumed by drift and
diffusion coefficients --- the raw $v_i$ may go slightly negative, but the
log-Euler price update preserves $S_{i+1} > 0$ exactly.

To make the simulator differentiable in $\theta = (\mu, \kappa, \theta_{\!H},
\xi, \rho, v_0)$, parameters are passed as PyTorch tensors with
\texttt{requires\_grad=True}. The only autograd-unfriendly operation in the
forward path is $\sqrt{v^+}$ at $v = 0$, which has infinite derivative; we
replace it with $\sqrt{v^+ + \varepsilon}$ with $\varepsilon = 10^{-12}$. The
numerical perturbation is well below the Monte Carlo noise floor.

\subsection{Energy-distance loss}

Given an observed log-return path $\{r_t\}_{t=1}^{T}$ and simulated returns
$\{\tilde r_j(\theta)\}_{j=1}^{M}$ from $M$ paths under candidate parameters
$\theta$, we minimize the energy distance
\begin{equation}
  \mathcal{L}_{\mathrm{ED}}(\theta)
  = \frac{2}{TM}\sum_{t,j} \lvert r_t - \tilde r_j(\theta) \rvert
  - \frac{1}{T^2}\sum_{t,t'} \lvert r_t - r_{t'} \rvert
  - \frac{1}{M^2}\sum_{j,j'} \lvert \tilde r_j(\theta) - \tilde r_{j'}(\theta) \rvert.
  \label{eq:energy}
\end{equation}
Equation~\eqref{eq:energy} is the squared MMD with the negative-Euclidean
kernel; it is zero iff the two distributions match \cite{szekely2013energy,
gretton2012mmd}. Compared to NLL, energy distance is more forgiving under
model misspecification --- mismatches in the deep tail contribute a bounded
amount rather than the unbounded penalty of KL divergence.

\subsection{Reparameterization and optimization}

Positive parameters are reparameterized via $\exp$, and $\rho \in (-1, 1)$
via $\tanh$. The optimizer (Adam) operates on unconstrained leaf tensors
$(\mu, \log\kappa, \log\theta_{\!H}, \log\xi, \log v_0, \tanh^{-1}\rho)$. To
bound the per-step cost of \eqref{eq:energy}, we subsample simulated returns
to a cap of $M = 2048$ before forming the pairwise-distance matrix.

\subsection{Why BPTT, not the stochastic adjoint}

For SDEs whose parameter dimension is large or whose horizon is long, the
stochastic adjoint of \cite{li2020adjointsde} reduces memory from
$O(N \lvert\theta\rvert)$ to $O(\lvert\theta\rvert)$ at the cost of a second
forward sweep. For Heston with $\lvert\theta\rvert = 6$ and $N \sim 250$
daily steps, the activation memory is below 1\,MB, so backprop-through-time
runs at roughly 1$\times$ cost rather than the adjoint's 2$\times$. We
therefore default to BPTT, with a finite-difference gradient-check script
verifying agreement with autograd to within Monte Carlo noise.

\section{VaR Forecasting and Backtesting}
\label{sec:var}

\subsection{Forecasters}

For each rolling window of length $W$ ending at time $t$, we form three
1-day-ahead VaR forecasts at level $\alpha$:
\begin{itemize}
\item \textbf{GBM-MLE.} Closed-form: $\hat\mu = \bar r / \Delta t +
  \tfrac{1}{2}\hat\sigma^2$, $\hat\sigma^2 = s_r^2 / \Delta t$, then
  $\widehat{\mathrm{VaR}}^{\mathrm{GBM}} = \hat\sigma\sqrt{\Delta t}\,
  \Phi^{-1}(\alpha) - (\hat\mu - \tfrac{1}{2}\hat\sigma^2)\Delta t$.
\item \textbf{Heston-AI.} Calibrate \eqref{eq:heston-S}--\eqref{eq:heston-rho}
  on $\{r_{t-W+1:t}\}$ by Adam through \eqref{eq:energy} (warm-started from
  the previous window). Simulate $M_{\mathrm{MC}} = 2 \times 10^4$ one-step
  paths under the fitted parameters and report the empirical
  $(1-\alpha)$-quantile of the simulated log-returns.
\item \textbf{Historical.} Empirical $(1-\alpha)$-quantile of the past $W$
  log-returns.
\end{itemize}

\subsection{Coverage tests}

A correctly specified VaR satisfies (i) unconditional coverage
$\mathbb{E}[I_t] = 1 - \alpha$ where $I_t = \mathbf{1}\{-r_t > \widehat{\mathrm{VaR}}_t\}$
and (ii) serial independence of $\{I_t\}$. We test (i) by the Kupiec POF
likelihood ratio and (ii) by the Christoffersen Markov-chain LR; their sum
gives the joint conditional-coverage statistic distributed $\chi^2_2$.

\section{Experiments}
\label{sec:experiments}

\subsection{Setup}

Data: daily adjusted close prices for SPY, AAPL, MSFT from 2018-01-01 through
2024-12-31, sourced via \texttt{yfinance}. Estimation window $W = 504$
trading days ($\sim$2 years), forecast cadence step = 21 days
($\sim$monthly), confidence level $\alpha = 0.95$.

Calibration hyperparameters: Adam learning rate $10^{-2}$, $n_{\mathrm{paths}}
= 256$ for the loss evaluation, 100 iterations per refit. Simulation seed
fixed for reproducibility; warm-start across windows.

\subsection{Calibration trajectories}

\textit{[Placeholder: figure
\texttt{figures/param\_trajectories\_SPY.pdf} produced by
\texttt{scripts/run\_backtest.py} showing $(\kappa, \theta, \xi, \rho, v_0)$
across rolling windows. Discussion: $(\kappa, \theta, v_0)$ are stable;
$(\rho, \xi)$ are noisier --- expected per the identifiability analysis in
\S\ref{sec:method}.]}

\subsection{Coverage results}

\textit{[Placeholder: table
\texttt{results/backtest\_alpha95.json} -> Table~\ref{tab:coverage}.]}

\begin{table}[h]
\centering
\caption{Out-of-sample VaR coverage at $\alpha = 0.95$. Bold p-values fail to
reject conditional-coverage at the 5\% level (i.e., the model passes).}
\label{tab:coverage}
\begin{tabular}{llrrrrr}
\toprule
Ticker & Method & Forecasts & Violations & Obs.\ rate & Kupiec p & CC p \\
\midrule
SPY    & GBM-MLE     & --- & --- & --- & --- & --- \\
SPY    & Heston-AI   & --- & --- & --- & --- & --- \\
SPY    & Historical  & --- & --- & --- & --- & --- \\
\addlinespace
AAPL   & GBM-MLE     & --- & --- & --- & --- & --- \\
AAPL   & Heston-AI   & --- & --- & --- & --- & --- \\
AAPL   & Historical  & --- & --- & --- & --- & --- \\
\addlinespace
MSFT   & GBM-MLE     & --- & --- & --- & --- & --- \\
MSFT   & Heston-AI   & --- & --- & --- & --- & --- \\
MSFT   & Historical  & --- & --- & --- & --- & --- \\
\bottomrule
\end{tabular}
\end{table}

\section{Results}
\label{sec:results}

\textit{[Discussion seeds.]}

\paragraph{Coverage.} Compare observed-rate columns to the nominal 5\%.
Heston-AI is expected to be better-calibrated than GBM-MLE on tickers with
visible volatility clustering (AAPL, MSFT through 2020--2022); historical
VaR's coverage tracks the unconditional empirical quantile and tends to
under-respond to recent vol shifts.

\paragraph{Independence.} Volatility-cluster periods (March 2020, late 2022)
should produce serially correlated breaches under GBM-MLE. The
Christoffersen test is the cleanest way to detect this; we expect Heston-AI
to pass independence where GBM-MLE fails.

\paragraph{Parameter stability.} The rolling-window plots demonstrate that
$\hat\rho$ and $\hat\xi$ have wider posteriors than $\hat\kappa, \hat\theta,
\hat v_0$, confirming the identifiability constraint discussed in
\S\ref{sec:method}. Joint calibration with implied-volatility surfaces (out
of scope here) is the standard fix.

\section{Discussion and Future Work}
\label{sec:discussion}

\paragraph{Limitations.} (i) Single-regime: Heston cannot represent regime
breaks (COVID March 2020) and absorbs them into inflated $\xi$.
(ii) Returns-only calibration: implied-vol data would tighten $(\rho, \xi)$
considerably. (iii) 1-day horizon: longer horizons compound the BPTT
activation memory; the stochastic adjoint becomes attractive.

\paragraph{Future work.} A C++/CUDA path-generation kernel (deferred to phase
2) would push the per-iteration cost down by an order of magnitude, opening
the door to higher-frequency calibration and joint multi-asset fits. A
Heston-with-jumps extension \cite{barndorff2002jumps} (Barndorff-Nielsen \&
Shephard) would directly address the regime-break failure mode.

\bibliographystyle{plain}
\bibliography{refs}

\end{document}
